\chapter{Healthy Food Swaps}

\section*{Definition and Overview}
Healthy food swaps are simple substitutions that replace less nutritious foods and drinks with healthier alternatives, without making big changes to your lifestyle. Swapping white bread for wholegrain, sugary drinks for water, and salty snacks for nuts or fruit are easy ways to improve the quality of your diet. Small, consistent swaps add up to meaningful improvements in nutrition, weight, and disease risk over time. The principle is to make gradual, sustainable changes that fit your tastes and budget, rather than overhauling your whole diet at once.

\section*{Epidemiology}
Most Australians eat more discretionary foods and less of the core food groups than recommended. Discretionary foods and drinks account for around a third of the energy intake of Australian adults, and sugary drinks are a major source of added sugar. Reducing these and replacing them with healthier options has a large potential to improve diet quality, reduce obesity, and lower the risk of heart disease and diabetes. Small swaps are an evidence-based strategy recommended by dietitians.

\section*{Aetiology and Pathophysiology}
Food swaps improve health through changes in nutrients and energy.

\begin{itemize}
    \item \textbf{Replacing added sugar:} swapping sugary drinks for water reduces sugar and energy intake
    \item \textbf{Replacing refined grains:} wholegrain swaps add fibre, which supports gut health and satiety
    \item \textbf{Replacing saturated fat:} choosing unsaturated fats supports heart health
    \item \textbf{Replacing salt:} lower-salt options reduce blood pressure
    \item \textbf{Energy balance:} lower-energy swaps support a healthy weight
    \item \textbf{Increasing nutrients:} swaps toward vegetables, fruit, and protein increase vitamins and minerals
    \item \textbf{Sustainability:} swaps that fit preferences are more likely to be maintained
\end{itemize}

\section*{Clinical Features}
Common situations where swaps are helpful include:

\begin{itemize}
    \item Replacing soft drink, cordial, and energy drinks with water
    \item Swapping white bread, rice, and pasta for wholegrain versions
    \item Replacing processed snacks with fruit, nuts, or yoghurt
    \item Choosing lean meats and reducing processed meats
    \item Reducing salt with herbs, spices, and lower-salt products
    \item Swapping full-cream dairy for reduced-fat where appropriate
    \item Choosing baked or grilled foods over fried foods
    \item Reducing takeaway and fast food
\end{itemize}

\section*{Differential Diagnosis}
Choosing the right swaps requires care.

\begin{itemize}
    \item "Healthy" marketed products that still contain high sugar or salt
    \item Portion size, which matters as much as the swap itself
    \item Individual needs, including diabetes, allergies, and dietary restrictions
    \item Budget and access, which affect the practicality of swaps
    \item Cultural and personal food preferences
    \item Swapping within the same food group (for example, wholegrain for white bread)
    \item The overall pattern of the diet, not just individual foods
\end{itemize}

\section*{When to Seek Medical Care}
Food swaps are self-directed for most people. See a doctor or dietitian for advice if you have a medical condition affected by diet, such as diabetes, kidney disease, or high cholesterol.

\textbf{Call 000 (or local emergency services) for:}
\begin{itemize}
    \item This chapter concerns nutrition rather than emergencies
    \item Severe allergic reactions to food require emergency care
    \item Any medical emergency unrelated to diet should be treated as usual
\end{itemize}

\section*{Diagnosis and Investigations}
Assessment helps target the swaps that matter most.

\begin{itemize}
    \item \textbf{Dietary review:} identifying the main sources of sugar, salt, and saturated fat in your diet
    \item \textbf{Food diary:} tracking intake for a few days reveals opportunities for swaps
    \item \textbf{Label reading:} comparing nutrition panels and ingredients
    \item \textbf{Medical review:} for conditions affected by diet
    \item \textbf{Dietitian advice:} individualised, practical suggestions
    \item \textbf{Follow-up:} reviewing progress and adjusting
\end{itemize}

\section*{Management}
Making effective swaps is about small, practical changes.

\begin{itemize}
    \item \textbf{Start with drinks:} replace sugary drinks with water, sparkling water, or milk
    \item \textbf{Upgrade grains:} choose wholegrain bread, rice, pasta, and cereals
    \item \textbf{Swap snacks:} replace chips and biscuits with fruit, nuts, yoghurt, or wholegrain crackers
    \item \textbf{Choose better protein:} lean meat, fish, eggs, legumes, and skinless poultry
    \item \textbf{Cut salt:} use herbs, spices, and lemon, and choose reduced-salt products
    \item \textbf{Cook healthier:} grill, bake, or steam instead of frying
    \item \textbf{Read labels:} check sugar, salt, and fat on packaged foods
    \item \textbf{Keep it sustainable:} choose swaps you actually enjoy
    \item \textbf{Gradual change:} introduce swaps one at a time
\end{itemize}

\section*{Complications}
\begin{itemize}
    \item Replacing one unhealthy food with another marketed as healthy
    \item Larger portions negating the benefit of a healthier choice
    \item Cost concerns if swaps are more expensive
    \item Nutrient imbalance from poorly planned swaps
    \item Giving up when swaps do not suit preferences
    \item Misleading "light" or "diet" products
\end{itemize}

\section*{Prognosis and Outcomes}
Healthy swaps produce measurable improvements in diet quality, weight, blood pressure, cholesterol, and blood glucose when maintained over time. Even small changes, such as replacing one sugary drink a day with water, can support weight management and reduce disease risk. Because swaps are gradual and flexible, they are easier to sustain than strict diets, and they build confidence for further change. Over years, consistent swaps contribute to better health and longevity.

\section*{Prevention and Self-Care}
\begin{itemize}
    \item Swap sugary drinks for water most of the time
    \item Choose wholegrain versions of bread, rice, and pasta
    \item Snack on fruit, nuts, or yoghurt instead of chips and biscuits
    \item Reduce processed and cured meats
    \item Season with herbs and spices instead of salt
    \item Grill, bake, or steam instead of frying
    \item Read food labels and compare products
    \item Make one or two swaps at a time
    \item Keep healthy options visible and easy to grab
    \item Enjoy treats occasionally without guilt
\end{itemize}

\section*{Sources and Further Reading}
The following reputable sources provide further information for healthcare students and patients seeking reliable, current advice about healthy food swaps.

\begin{itemize}
    \item Dietitians Australia. \textit{Healthy food swaps.} https://dietitiansaustralia.org.au/
    \item Eat for Health (Australian Government). \textit{Healthy eating advice.} https://www.eatforhealth.gov.au/
    \item Healthdirect Australia. \textit{Healthy food swaps.} https://www.healthdirect.gov.au/
    \item Australian Government Department of Health and Aged Care. \textit{Healthy eating and nutrition.} https://www.health.gov.au/
    \item NHS. \textit{Healthy food swaps.} https://www.nhs.uk/
    \item American Heart Association. \textit{Healthy substitutions.} https://www.heart.org/
\end{itemize}
